% Chapter 1: Why numerical derivatives (fundamentals author). PLAN.md section 1.4.
% Uses F1 (fig:secant), S1 (sb:alternatives), E1 (code/01-why-first.cpp and code/shomate.hpp), T0 (tab:method-summary).
% Language pass for students (PLAN.md section 12.9): the wording only; content, evidence and order as at 3192e8a.
\chapter{Why numerical derivatives}\label{ch:why}

Engineering models are full of derivatives that nobody writes down. A heat capacity is the slope of an enthalpy. A
Jacobian is the slope of a set of balance equations.\mnote{A \emph{Jacobian} is the matrix of the partial derivatives
of a set of equations with respect to their unknowns.} A sensitivity is the slope of a result with respect to one of
the parameters that went into it. When the function is a formula on paper, its derivative is a matter of calculus.
When the function is a computer program, it often is not.

This chapter first shows where such derivatives arise. It then compares the tools that can compute them, and
describes what the library \nxx{} offers. It ends with a first computation and a half-page summary for the reader in
a hurry.

\section{Derivatives without a formula}\label{sec:why-without-formula}

In engineering software, a function is usually a piece of code: it takes a number and returns a number. Here are
some typical cases.
\begin{itemize}
  \item \emph{Property models.} A heat capacity $\cp = \dd \enth / \dd T$ from an enthalpy $\enth(T)$; a
    compressibility from an equation of state; the slope of a vapour pressure for the Clausius--Clapeyron equation.
    Such a model is a correlation, a formula fitted to measured data, inside a property library. It often comes in
    several pieces.
  \item \emph{Simulators.} The outlet temperature of a heat exchanger, as a function of a feed flow. Each value
    comes from a simulation that iterates until it converges.
  \item \emph{Implicit functions.} A friction factor is defined by an equation that it must satisfy. A flash
    temperature is defined by an energy balance. In both cases the value comes from a solver, and no formula for it
    exists.
  \item \emph{Sensitivities.} How much does a reaction rate change per kelvin, or per unit of activation energy? A
    design review asks for numbers like these.
  \item \emph{Newton's method.} To solve $g(T) = 0$, Newton's method needs the slope $g'(T)$ at every iteration.
\end{itemize}
In every case we can \emph{evaluate} the function at any point we like, but nobody gives us its derivative. We could
work the derivative out by hand from the code. That is tedious and error-prone, and impossible when the code is not
ours.

The way out is old and simple. The derivative $f'(x)$ is the slope of the tangent to the curve of $f$ at $x$. A line
through two nearby points of the curve is called a \emph{secant}, and its slope is close to that of the tangent
(\cref{fig:secant}).\mnote{The \emph{step} $h$ is the distance between $x$ and the points at which $f$ is
evaluated.} So we evaluate $f$ at $x$ and at $x + h$, and divide the change of $f$ by $h$. The result is an estimate
of $f'(x)$, called a \emph{finite difference}. It is the difference quotient from the definition of the derivative,
taken at a small but finite $h$ instead of in the limit $h \to 0$.

The figure shows a second secant, through $x - h$ and $x + h$, placed symmetrically around $x$. It is much closer to
the tangent than the secant through $x$ and $x + h$. \Cref{ch:taylor} explains why. Since the derivative is a limit,
it is tempting to make the step as small as we like. \Cref{ch:rounding} explains why we cannot simply do that.

% A text-width figure in the body, not a margin figure (the author's release comment: too small to read there).
\begin{figure}
\sidecaption{The tangent at $x$ and two secants. The central secant, through $x - h$ and $x + h$, is nearly parallel
to the tangent. The forward secant, through $x$ and $x + h$, is visibly steeper. A diagram, without
data.\label{fig:secant}}
\centering
% F1 fig:secant (Chapter 1, margin): the tangent at x and two secants, a diagram without data. The curve is
% exp(t) around x = 0 with a step of 0.7; any smooth curve would do, and the text does not name it. The forward
% secant joins the points at x and x + h; the central secant joins the points at x - h and x + h and is nearly
% parallel to the tangent: the one thing to notice, in accent. The lines end at staggered abscissas so that their
% labels do not collide.
\begin{tikzpicture}[x=17mm, y=10.5mm, font=\sffamily\scriptsize]
  \pgfmathsetmacro{\stepH}{0.7}
  \pgfmathsetmacro{\fMinus}{exp(-\stepH)}
  \pgfmathsetmacro{\fPlus}{exp(\stepH)}
  \pgfmathsetmacro{\slopeF}{(\fPlus - 1) / \stepH}              % slope of the forward secant
  \pgfmathsetmacro{\slopeC}{(\fPlus - \fMinus) / (2 * \stepH)}  % slope of the central secant
  % the axis with the three nodes
  \draw[thin, black!55] (-1.2, -0.3) -- (1.25, -0.3);
  \foreach \nodeX/\nodeLabel in {-\stepH/{$x - h$}, 0/{$x$}, \stepH/{$x + h$}} {
    \draw[thin, black!55] (\nodeX, -0.3) -- ++(0, -0.05)
      node[below, black!70, inner sep=1.5pt] {\nodeLabel};
    \draw[densely dotted, black!40] (\nodeX, -0.3) -- (\nodeX, {exp(\nodeX)});
  }
  % the function
  \draw[line width=0.8pt, black!80, domain=-1.2:1.22, samples=60, smooth] plot (\x, {exp(\x)})
    node[left, black!80, inner sep=2pt] {$f$};
  % the tangent at x: slope f'(x) = 1
  \draw[line width=0.6pt, black!80] (-1.15, {1 - 1.15}) -- (0.95, {1 + 0.95})
    node[right, inner sep=1.5pt] {tangent};
  % the central secant through x - h and x + h
  \draw[line width=0.6pt, accent] (-1.15, {\fMinus + (-1.15 + \stepH) * \slopeC})
    -- (1.05, {\fMinus + (1.05 + \stepH) * \slopeC})
    node[right, inner sep=1.5pt] {central};
  % the forward secant through x and x + h
  \draw[line width=0.6pt, muted] (-0.3, {1 - 0.3 * \slopeF}) -- (1.15, {1 + 1.15 * \slopeF})
    node[right, inner sep=1.5pt] {forward};
  % the sampled points
  \fill[muted] (\stepH, \fPlus) circle (1.1pt);
  \fill[accent] (-\stepH, \fMinus) circle (1.1pt);
  \fill[black!80] (0, 1) circle (1.1pt);
\end{tikzpicture}

\end{figure}

\section{The alternatives in one page}\label{sec:why-alternatives}

A finite difference is not the only way to compute a derivative, and it is not always the best. This section
compares it with three alternatives, so that you can tell which tool fits which case.
\begin{itemize}
  \item \emph{The analytic derivative} is the formula for $f'$, found by hand or by computer algebra. It is exact and
    cheap to evaluate. Use it when the formula for $f$ is at hand and short enough to differentiate without
    mistakes.
  \item \emph{Automatic differentiation} runs the program with a special number type. Each number of this type
    carries a derivative along with its value. The result is exact to rounding.\mnote{\emph{Exact to rounding}: the
    only error is the rounding that the computer applies to every arithmetic result, because it keeps a fixed number
    of digits.} The method needs the source code of the program, written for a generic number type, that is, so
    that it works with any number type.
  \item \emph{The complex step} evaluates $f$ once, at a complex argument, and reads $f'$ from the imaginary part of
    the result. It too is exact to rounding. It needs an $f$ that accepts complex numbers and is analytic, that is,
    differentiable as a function of a complex variable.
  \item \emph{A finite difference} needs only values of $f$ at real points. So it works on any function that the
    program can evaluate: a library call, a solver, legacy code. It has two prices. It is less accurate: it keeps
    about two thirds of the digits of the number type\mnote{A \emph{number type} such as \code{double} stores a real
    number with a fixed number of significant digits.} (\cref{ch:rounding}). And it needs two or more evaluations of
    $f$.
\end{itemize}
This book is about the last of the four. It is the tool for a function that is a \emph{black box}: we can call it,
but we cannot look inside. When one of the first three is open to you, it is better. The sidebar below says when
(\cref{sb:alternatives}).

\begin{sidebar}[label=sb:alternatives]{When automatic differentiation or the complex step is the better tool}
\emph{Forward-mode automatic differentiation} replaces each number $v$ in the computation by a pair $(v, v')$: the
value and its derivative. It replaces each operation by its rule of calculus. For a product the rule is
$(u, u') \cdot (v, v') = (uv, u'v + uv')$, and so on for the other operations. We seed the input $x$ with the pair
$(x, 1)$, since the derivative of $x$ with respect to itself is 1. At the end the computation delivers the pair
$(f(x), f'(x))$. The derivative is exact to rounding, and it costs a small multiple of one evaluation. The method
needs the function's source, written for a generic number type (a C++ template, for instance), and a library that
supplies the pair type.\cite{griewank2008}

The \emph{complex step} uses $f'(x) \approx \operatorname{Im} f(x + \mathrm{i} h) / h$.\cite{squire1998} The formula
subtracts no two nearly equal values. So the step can be taken far below the limit that rounding sets for a finite
difference, and one evaluation gives a result that is exact to rounding. It needs $f$ implemented for complex
arguments, and analytic near $x$. An absolute value, a comparison or a table lookup inside $f$ breaks it. The complex
step is closely related to forward-mode automatic differentiation.\cite{martins2003} It is also a special case of the
earlier contour-integral methods of Lyness and Moler, which give higher derivatives of analytic functions as
well.\cite{lyness1967} Squire and Trapp brought the simple first-derivative form back into use.

\code{nxx::deriv} takes any callable \code{fn}, such as a function or a lambda, that maps a real number to a real
number. It assumes nothing else about \code{fn}: not its source, not a complex extension, not even smoothness beyond
what the error terms of the later chapters need. That is the case that the other tools do not cover.
\end{sidebar}

\section{What the library offers}\label{sec:why-offers}

\code{nxx::deriv} is the numerical-derivatives module of \nxx. \nxx{} is a header-only C++23
library.\mnote{\emph{Header-only}: the library is a set of header files that your program includes; there is nothing
to compile or link separately.} Everything in the module goes through one function, \code{derivative}. These are the
choices it offers.
\begin{itemize}
  \item \emph{Six methods.} \code{Central<P>}, the default, samples $f$ on both sides of $x$. \code{Forward<P>} and
    \code{Backward<P>} sample on one side only; they are for a function that is defined on one side of $x$.
    \code{Ridders} adapts its step and also estimates its own error. Its one-sided forms, \code{ForwardRidders} and
    \code{BackwardRidders}, do the same on one side of $x$ (\cref{sec:ridders-onesided}). The first three methods
    are \emph{stencils}.\mnote{A \emph{stencil} samples $f$ at a fixed pattern of points around $x$ and combines
    the samples with fixed weights. Ridders' method decides its points as it goes.} The \code{P} in their names is
    the accuracy order, which \cref{ch:taylor} defines.
  \item \emph{Derivatives of order 1 to 4}, $f'$ to $f^{(4)}$.
  \item \emph{Any real type}: \code{float}, \code{double}, or a user-defined type with more digits, such as those
    of Boost.Multiprecision.
  \item \emph{No exceptions from the library.} A failure is returned as a value. The result is a
    \code{std::expected}, which holds either an estimate or the error \code{non_finite}. The library throws nothing
    itself. An exception thrown by \code{fn} passes through unchanged, and \code{derivative} is not
    \code{noexcept}.\mnote{A C++ function marked \code{noexcept} promises never to throw an exception.}
  \item \emph{A default step}, so that a call needs only the function and the point.
\end{itemize}
\Cref{ch:using} shows how to make the calls, and \cref{ch:api} lists every one of them.

\section{A first example}\label{sec:why-first}

One example runs through all of Part~I: the heat capacity of a gas, computed from its enthalpy. We call it the
nitrogen example. For nitrogen, the NIST Chemistry WebBook gives the molar enthalpy in the Shomate
form.\mnote{The \emph{Shomate form}: a fitted polynomial in $T/1000$ plus a $1/T$ term, one set of coefficients per
temperature range.} The WebBook's formula is a fit to the NIST-JANAF thermochemical
tables.\cite{nist-webbook,chase1998} We use the fit for the range \SIrange{500}{2000}{\kelvin} and the point
$T_0 = \SI{1000}{\kelvin}$. The function returns $\enth(T)$ in \si{\joule\per\mole}, as the enthalpy relative to its
value at \SI{298.15}{\kelvin}. Then $\dd \enth / \dd T$ is $\cp$ in \si{\joule\per\mole\per\kelvin} directly, with no
conversion. The whole program is short:
\snippet{01-why-program}{program}
It prints $\cp = \num{\valCpDefaultValueLong}$\,\si{\joule\per\mole\per\kelvin}. The heat capacity takes one call.
We gave \code{derivative} no method and no step, so it uses its defaults. The method is \code{Central<2>}, the
central difference of \cref{fig:secant}. The step is the default step, here \SI{\valCpDefaultStep}{\kelvin};
\cref{ch:rounding} explains where it comes from.

The result \code{d} holds either an estimate or the reason why there is none. The test \code{if (d)} is true when
there is an estimate. Only then does the program read \code{d->value}, the derivative. The \code{else} branch handles
the failure. The test is not a formality, because a call can fail. It then returns the error
\code{error::non_finite}. This happens, for example, when \code{fn} returns an infinite or NaN
value\mnote{\emph{NaN}, ``not a number'', is the value of an undefined floating-point result such as $0/0$.} at a
point that the method samples. It also happens when the step is too small to move $x$ (\cref{ch:rounding}). Every
listing in Part~I reads a result this way.

The first header that the program includes is the book's own, not the library's. \code{book::} is the namespace of
the book's examples. The header \code{shomate.hpp} defines the fit once for all of them. It holds the WebBook's
coefficients and its expression for $\enth$. First the coefficients:
\snippetheader{shomate}{shomate}
\needspace{4\baselineskip}% the lead-in with the first lines of the listing, so no empty frame top ends a page
Then the enthalpy itself, in \si{\joule\per\mole}:
\snippetheader{shomate}{shomate-H}

How good is the estimate? For this function we know the answer. The WebBook also publishes the fit's $\cp(T)$,
which is the exact derivative of $\enth(T)$. It gives $\cp(T_0) = \num{\valCpExactLong}$\,\si{\joule\per\mole\per\kelvin}.
The estimate is off by \num{\valCpDefaultRelError} relative.\mnote{A \emph{relative} error is the error divided by
the size of the exact value; it has no unit.} This closed form is the \emph{answer key} of the book. We differentiate
a function whose derivative we know, so that we can measure every error. In practice you would differentiate one
whose derivative you do not know. \Cref{ch:using} does so, with a friction factor that comes from a solver.

Note that the closed form is exact for the \emph{function}, not for nitrogen. At \SI{1000}{\kelvin} the tables give
\SI{\valCpTableValue}{\joule\per\mole\per\kelvin} (Part~II, p.~1621),\cite{chase1998} and the fit differs from this
by \num{\valCpFitVsTableRelDiff} relative. That difference is a matter for the property data, not for the derivative.

Now we try Ridders' method. It adapts its step and reports an estimate of its own error in \code{abs_error}:
\snippet{01-why-first}{first-ridders}
The method evaluates $\enth$ \num{\valCpRiddersEvals} times, where the stencil needed two. It returns
$\cp = \num{\valCpRiddersValue}$\,\si{\joule\per\mole\per\kelvin} with an error estimate\mnote{\code{abs_error} is a
\code{std::optional}: it either holds a value or is empty. It is empty for the stencils. The \code{*} in
\code{*r->abs_error} reads the value.} of \SI{\valCpRiddersEstimate}{\joule\per\mole\per\kelvin}. The actual error is
\SI{\valCpRiddersActual}{\joule\per\mole\per\kelvin}, or \num{\valCpRiddersRelError} relative. So the estimate is
honest here: it is larger than the actual error. \Cref{ch:ridders} shows how the method works, and how far we can
trust its estimate.

\section{If you read nothing else}\label{sec:why-summary}

This section gives the essentials of the book in half a page: which method to use, what accuracy and cost to expect,
and how to check a result. \Cref{tab:method-summary} chooses a method for the five situations that cover most uses. The full
decision table is \cref{tab:method-choice} in \cref{ch:using}.

\paragraph{Accuracy.} At the default step, a first derivative in \code{double} has about
\num{\valHeadlineFirstDigits} correct digits,\mnote{\emph{About $n$ correct digits}: a relative error of about
$10^{-n}$.} and a fourth derivative about \num{\valHeadlineFourthDigits}. Both are the medians of
\cref{tab:accuracy-matrix}, measured on $\mathrm{e}^x$. Higher derivatives lose digits fast (\cref{ch:orders}).

\paragraph{Cost.} We count cost in evaluations, that is, in calls of $f$. A stencil costs a fixed number of
evaluations: two for \code{Central<2>} and $f'$. Ridders' method typically needs \num{\valRiddersEvalsMedianFirst}
for $f'$, and never more than \num{\valRiddersMaxEvalsFirst} (\cref{tab:cost}).

\paragraph{Did it work?} Test the result before you use it, as the program of \cref{sec:why-first} does. Write
\code{if (d)}, read \code{d->value} only inside it, and handle the failure in an \code{else}.

\paragraph{Is it right?} A value can be wrong without any warning. Three situations cause most such errors:
\begin{itemize}
  \item $f$ changes appreciably over a distance in $x$ that is far from $\max(\abs{x}, 1)$ in your units, as across
    a narrow transition or for a small mole fraction (\cref{sec:rounding-scale});
  \item $f$ is the output of a solver, or is otherwise less accurate than the number type
    (\cref{sec:rounding-noise});
  \item $f$ has a kink or a range break within a few steps of $x$ (\cref{sec:using-pitfalls}).\mnote{A \emph{kink}
    is a point where $f$ is continuous but its slope jumps. A \emph{range break} is the point where one fitted formula
    hands over to the next.}
\end{itemize}
How can you tell? A stencil reports no error estimate. If in doubt, repeat the call with the steps $2h$ and $h/2$,
and count the digits that agree (\cref{sec:using-check}). Ridders' method does report an estimate: read it before you
report a value. Accept the result \code{res} only if
\code{res && res->abs_error && *res->abs_error <= atol + rtol * std::abs(res->value)} (\cref{sec:using-result}).
Here \code{rtol} is a relative tolerance and \code{atol} an absolute one. Take \code{rtol} $= 10^{-n}$ for a value
that you report to $n$ significant figures. Take \code{atol} as the error you can accept, in the derivative's units.
Near a zero of $f'$ the relative part of the tolerance shrinks to almost nothing, so \code{atol} is what decides
there. Inside an iteration, no such check is needed. Keep in mind what a passed check says: that the estimate is
small. It does not say that the value is right where the estimate is blind, as at a kink.

% The table follows the two checks so that the list of the three silent failures is never split by
% the float (R2-engineer-03); it still precedes the paragraph that cites it on the page. [htb]: here, after the
% checks, rather than at the foot of the chapter's last page below a gap (the language pass, PLAN.md 12.9).
\begin{table}[htb]
\sidecaption{The method in five situations. Each row is a call \code{derivative<M>(fn, x)}, or
\code{derivative<M>(fn, x, h)} when the step is given. Here $\epsmach$ is the machine epsilon of the type, the
distance from 1 to the next number of the type ($2^{-52}$ for \code{double}). The scale $w$ is the distance in $x$
over which $f$ changes appreciably (\cref{sec:rounding-scale}). The noise level $\noise$ is the relative accuracy of
$f$, which says how many digits of it you can trust: for example $10^{-8}$ for a solver converged to a relative tolerance of $10^{-8}$ on its
output (\cref{sec:rounding-noise}). The evaluations are those of \code{fn} for $f'$. The full table, with the
evidence for each row, is \cref{tab:method-choice}. The quick-reference page at the end of the book repeats this
one.\label{tab:method-summary}}
\centering\footnotesize
\renewcommand{\arraystretch}{1.15}
\begin{tabular}{@{}>{\raggedright\arraybackslash}p{33mm}>{\raggedright\arraybackslash}p{36mm}%
  >{\raggedright\arraybackslash}p{36mm}@{}}
\toprule
Situation & Method \code{M} and step & Cost and accuracy\\
\midrule
\tablerows{chapters/shared-method-summary}
\bottomrule
\end{tabular}
\end{table}

\paragraph{Where to read on.} \Cref{ch:taylor} explains why the central difference beats the forward one.
\Cref{ch:rounding} explains why a smaller step is not always better, and when to pass your own. \Cref{ch:stencils}
adds more points and one-sided rules, and \cref{ch:orders} treats higher derivatives. \Cref{ch:ridders} covers
Ridders' method, and \cref{ch:using} the calls and the pitfalls.
